\documentclass[prep,both]{debate}
\motion{超级英雄的出现，对世界而言更是幸福 / 更是灾难}
\meta{赛制}{中国科学技术大学新生辩论赛（默认赛制，四对四）}
\meta{生成日期}{2026-09-15}
\begin{document}
\debatetitle
\begin{summary}
\begin{fields}
\claim[持方优劣评估]{\textbf{正方（更是幸福）较劣势}，劣势主要来自定义空间与论据可得性。常识里的超级英雄是“救人的好人”，正方起手有评委直觉红利；但中立定义里“不受既有授权与问责”这一条是常识的一部分（所有人都知道蝙蝠侠不归警察管），这一条一旦写进定义，反方的战场就自动打开，而且现实世界里“不受约束的执法者”的经验证据极其丰富，正方几乎找不到“超常能力带来净收益”的直接对应物。正方的打法是：正一自己把这一条说出来并立刻用判准框住它，全场只打“最坏的那一类伤害”与“时间窗口”，反复追问反方“你说的灾难是和哪个世界比”。}
\thesis{正方}{跟没有超级英雄的世界比，他挡下的是那些本来必然发生、而且谁都来不及阻止的伤害，这份净改善大过他带来的秩序代价。 }
\thesis{反方}{超级英雄的出现，第一次把一种压倒性的力量放在了没有任何人能够纠错的位置上，世界从此失去的不是安全感，是出错以后改回来的能力。 }
\end{fields}
\end{summary}
\section{一、赛制要点}
\begin{flowtable}
\block{A}
\flow{1 正一立论 → 2 反四质询正一 → 3 反一立论 → 4 正四质询反一}{180 / 120 / 180 / 120 秒}{正一、反四、反一、正四}{立论清楚、判准站得住；质询拿到对方至少一个关键承认。本题的关键承认只有两个：正方承认“超级英雄不受既有权力问责”，反方承认“没有超级英雄的世界里有些人确实救不下来”}
\block{B}
\flow{6 反二驳论 → 7 正二驳论 → 8 二辩对辩}{120 / 120 / 各 90 秒}{反二、正二}{反二打正方立论里最难圆的一点（谁来纠错）；正二先回应反二再打反方立论（你的灾难是和哪个世界比）}
\block{C}
\flow{10 正三盘问 → 11 反三盘问 → 12 正三小结 → 13 反三小结}{各 90 秒}{正三、反三}{盘问对方二辩与四辩同一个问题，抓口径不一致。本题最容易抓到的不一致：反方二辩说“他不一定作恶”、反方四辩说“他迟早作恶”}
\block{D}
\flow{15 自由辩论}{双方各 240 秒}{全员}{主战场推进到位，必问问完，对方必答问题正面回应过}
\block{E}
\flow{17 反四结辩 → 18 正四结辩}{各 210 秒}{反四、正四}{回到判准与一句话立论，处理掉对方全场最强的一个点}
\end{flowtable}
计时方式对策略的影响：质询是双边计时，被质询方拖时间就是在吃质询方的时间，所以本题的质询问题必须是封闭式短问，尤其是问“超级英雄受不受现有法律问责”这种一句话能答的问题，对方一旦展开就立刻打断。对辩是单边计时，比的是每次发言的密度，不是嗓门。反四先结辩，所以反四必须写封路式结辩：把正四最可能的收束方向（“跟现实比不是跟理想比”）提前点破。

\section{二、辩题性质与争议焦点}
\begin{fields}
\claim[辩题性质]{\textbf{价值比较题}（外壳是假设性的利弊比较）。“更是幸福 / 更是灾难”这个句式意味着双方都承认另一面存在，争的是哪一面更主导、更持久、更难逆转。任何一方试图证明“完全没有坏处”或“完全没有好处”，都是自己加重义务。}
\field{正方倾向把题目解释成}{把“有超级英雄的世界”和“没有超级英雄的世界”放在一起比的利弊题，衡量单位是可数的伤害与生命。
}
\field{反方倾向解释成}{对世界秩序的结构性影响评估，衡量单位是纠错能力与长期风险，看的是最坏情形而不是平均情形。
}
\field{争议焦点 （双方必然相遇的三个战场）}{
\begin{fields}
\field{谁来纠错}{一种不受既有授权与问责的力量，能不能是安全的。
}
\field{制度与人的存废}{他的存在让普通制度和普通人更强还是更弱。
}
\field{风险的不对称}{一个坏的持有者造成的损失，能不能被许多好的持有者的收益抵消。
}
\end{fields}
}
\end{fields}
\subsection{本题的证据纪律（全队必须遵守）}
这是一道假设性辩题，所以论据规矩比平时更严：

\begin{bullets}
\item \textbf{漫威、DC 的情节是设定，不是论据。} 剧情是编剧为了讲故事设计的，用“复仇者联盟打赢了灭霸”证明现实结论，等于用小说证明历史。评委听得出来，而且会扣分。
\item 漫画与电影\textbf{只在举例说明“我方说的是哪一类人”时使用}，每次使用都要明说“这是设定，不是我方论据”。
\item 诚实的论据基础只有现实世界的对应物，本文档的全部数据都取自这五类：\textbf{私刑与义警研究}、\textbf{核武器与不对称力量}、\textbf{强人政治与制度替代}、\textbf{榜样叙事对助人行为的实验}、\textbf{超级大国单边干预与不干预的后果}。
\item 一个听起来合理但没有出处的数字，在这道题上比没有数字更糟。对方只要问一句“出处是什么”，全队的可信度就一起掉。
\end{bullets}
\section{三、关键词定义}
\subsection{3.1 超级英雄}
\begin{fields}
\field{调研}{
\begin{bullets}
\item 词典义：《不列颠百科词典》（Britannica Dictionary）“superhero” 有两义。义一：“a fictional character who has amazing powers (such as the ability to fly)”（拥有惊人能力的虚构人物，例如会飞）。义二：“a very heroic person”（非常英勇的人），例句是“从事这项危险工作的人是现实生活中的超级英雄”。已核实 2026-09-15。
\item 义二是\textbf{比喻义}，两方都不能用：如果超级英雄就是“很英勇的人”，那消防员、医生已经是超级英雄，辩题变成“世上有勇敢的人是好事还是坏事”，没得打。
\item 学术义：本题在现实世界里最接近的学术概念是 \textbf{vigilantism（义警行为）}。Bateson (2021) 把它定义为“在法律程序之外对违法行为的预防、调查或惩罚”（the extralegal prevention, investigation, or punishment of offenses）。这个定义的价值在于：它不问动机好坏，只问行为是否在既有法律程序之外。有把握，见第十一节。
\end{bullets}
}
\begin{procard}{正方有利定义}\definition{超级英雄，指拥有超常能力、并且以利他为动机持续行善的人。}\end{procard}
\begin{bullets}
\item 有利之处：把“做好事”写进定义，灾难在定义上就很难成立。
\item \textbf{反方如何反驳}：这是定义杀。你把结论藏进了定义——一个“定义上永远做对事的人”出现，当然不会是灾难，那这场辩论不用打了。而且辩题问的是“超级英雄的出现”这件事本身对世界的意义，不是“一个好人出现”。请对方给出一个反方还能成立的定义。
\end{bullets}
\begin{concard}{反方有利定义}\definition{超级英雄，指不受任何既有权力约束、可以单方面对他人使用压倒性力量的个体。}\end{concard}
\begin{bullets}
\item 有利之处：把不可问责写进定义，灾难几乎自动成立。
\item \textbf{正方如何反驳}：你把“英雄”两个字删掉了，只留下“超能力者”。常识里的超级英雄是去救人的，不是去掌权的。按你的定义，一个从不出手的超能力者也叫超级英雄，这不是这道题在讨论的东西。
\end{bullets}
\begin{neutralcard}{中立定义（一辩稿用这一版）}\textbf{超级英雄，指拥有远超常人的能力，并且主动用这种能力介入公共安全事务的个体；他行动的依据是自己的判断，既不由现有的法律和权力机构授权，也不由它们问责。} 三条缺一不可：能力超常、主动介入、不受既有授权与问责。\end{neutralcard}
\begin{bullets}
\item 合理性：\textbf{贴近常识}——每个人脑子里的超级英雄都同时满足这三条，蝙蝠侠不归警察管是所有人都知道的事，把这一条删掉反而不像超级英雄了。\textbf{让双方都有得打}——正方可以论证他救了人，反方可以论证他不可控，谁都没被定义排除。\textbf{排除了比喻义}——消防员受制度授权也受制度问责，不在此列。
\end{bullets}
\end{fields}
\subsection{3.2 出现}
\begin{fields}
\field{调研}{辩题用“出现”而不是“出手”“救人”“存在”，说明讨论对象是一个\textbf{事实的确立}，不是一次事件。
}
\begin{procard}{正方有利定义}\definition{出现 = 在危难现场出现，强调具体的救助行为。}\end{procard}
\begin{bullets}
\item 有利之处：把辩论拉到“救没救到人”上，正方的可数收益最大化。
\item \textbf{反方如何反驳}：那只是“一次出手”。题目说的是“出现”，是这种人存在于世界上这件事本身，包括所有人因此改变的行为。按你的定义，超级英雄回家睡觉那天世界就不受影响了，这说不通。
\end{bullets}
\begin{concard}{反方有利定义}\definition{出现 = 超常能力被证明可以存在于个体身上，因此它也会出现在其他人身上。}\end{concard}
\begin{bullets}
\item 有利之处：直接推出扩散与军备竞赛。
\item \textbf{正方如何反驳}：辩题说的是“超级英雄的出现”，不是“超能力的普及”。你从一个人滑到所有人，这一步没有论证。
\end{bullets}
\begin{neutralcard}{中立定义}\textbf{出现，指这种人存在于世界上成为一个既定事实，并且被世界知道。} 判断对象是这个事实带来的全部后果，既包括他做了什么，也包括别人因为他而改变了什么。\end{neutralcard}
\begin{bullets}
\item 合理性：贴近“出现”二字的字面意思，也贴近双方真正想讨论的东西；它不预设他一定出手，也不预设能力会扩散。
\end{bullets}
\end{fields}
\subsection{3.3 世界、更是幸福、更是灾难}
\begin{fields}
\field{世界}{全人类，以及人类共同生活的秩序。不是某一次事件里被救下的那些人。
}
\field{幸福}{不是快乐的情绪，是世界整体处境变好——人更安全，也更能靠自己决定怎么活。
}
\field{灾难}{不是“存在风险”，是整体处境出现\textbf{重大且难以逆转}的恶化。
\begin{bullets}
\item 正方对反方定义的反驳：如果“只要存在任何风险就是灾难”，那汽车、电、互联网全都是灾难，这个词就没有意义了。
\item 反方对正方定义的反驳：如果“只要救了人就是幸福”，那你只数救下来的人，不数被改变的秩序，等于把一半账目撕掉。
\end{bullets}
}
\field{更是}{双方都承认两面并存，比的是哪一面更主导。谁都不能只证明“有好处”或“有坏处”。
}
\end{fields}
\subsection{3.4 隐含要素}
\begin{fields}
\field{主体}{整个世界与其中的普通人，不是被救的个别当事人。
}
\field{范围}{全球，不限国别。
}
\field{时间尺度}{从超级英雄出现那天起的长期。短期与长期背离时以长期为准，这一条对反方有利，正方要在判准里争。
}
\field{比较对象}{\textbf{没有超级英雄的现实世界}，不是一个能及时救人的理想制度。这是全场最容易被偷换的一点，正方要反复钉住，反方要反复论证“出现之后比出现之前更糟”，而不是“比理想世界糟”。
}
\end{fields}
\section{四、判准与论证义务}
\begin{fields}
\claim[判准是什么]{这道题是价值比较题，评委要判的不是“有没有好处坏处”，而是“净下来更像哪一种”。所以判准必须同时能装下收益和代价，并且给出一个比较的尺度。}
\begin{procard}{正方有利判准}\definition{看超级英雄出现后，那些本来必然发生的伤害有没有被挡下来，以可数的生命与损失为准。}\end{procard}
\begin{bullets}
\item 有利之处：收益可数，代价不可数，可数的东西评委更容易记分。
\item \textbf{反方如何反驳}：这个标准只数救下来的人，不数因此改变的秩序，等于规定了只看一半账。而且它把比较对象偷换成“什么都不做”，可现实世界里从来不是什么都不做。
\end{bullets}
\begin{concard}{反方有利判准}\definition{看世界有没有出现一种不可问责、不可制衡的力量，只要出现就是灾难。}\end{concard}
\begin{bullets}
\item 有利之处：一旦接受，反方只需要论证定义里已有的东西。
\item \textbf{正方如何反驳}：这是把“存在风险”直接等于“灾难”，取消了程度比较，而“更是”这个词要求的恰恰是程度比较。按这个标准，任何一项新能力刚出现时都是灾难。
\end{bullets}
\begin{neutralcard}{中立判准（一辩稿用这一版）}\textbf{把有超级英雄的世界和没有超级英雄的世界放在一起，长期看，普通人的处境是变好了还是变坏了。} 处境有两项：一是\textbf{安全}，能不能不被伤害；二是\textbf{自主}，能不能靠自己和自己的制度决定怎么活。两项加起来的净变化落在哪一边，哪一边赢。\end{neutralcard}
\begin{bullets}
\item 合理性：它同时装下了双方最在乎的东西（正方在乎安全，反方在乎自主）；它照顾的是题目里的主体“世界”而不是个别当事人；它能真正区分双方，因为双方对这两项的预测确实相反。
\item 在中立判准下，\textbf{正方需要证明}：超级英雄挡下的伤害是真实且重大的，并且他带来的秩序代价是可控的、没有大到抵消这份改善。只证明“救了人”不够。
\item 在中立判准下，\textbf{反方需要证明}：超级英雄带来的秩序代价是结构性的、难以逆转的，并且大到把救援带来的改善抵消掉甚至倒挂。只证明“有风险”不够。
\end{bullets}
\end{fields}
\prosection{五、正方论点}
\prosubsection{论点一：救得下来}
\begin{fields}
\claim[主张]{世界上有一类伤害，能不能停下来只取决于那个时间点有没有人具备相应的能力；超级英雄补的正是这个缺口。}
\begin{mechanism}[机制]
\step 这类伤害的瓶颈是能力本身，不是道理
\step 现有的技术、制度与人的身体极限，让这个能力缺口在每一个具体的时间窗口里都补不上
\step 超级英雄出现，这一类伤害就从“来不及”变成“来得及”。

\end{mechanism}
\begin{evidence}{学理}[有把握]
公共品供给不足与搭便车（曼瑟尔·奥尔森，1965，《集体行动的逻辑》）。\sub{核心主张}{人人受益的东西，人人都指望别人去做，所以长期供给不足。}\sub{支撑}{解释了为什么“能力缺口”不是偶然的疏忽，而是一个结构性的空洞。}\sub{边界}{奥尔森讲的是意愿问题，不是能力问题，反方会用这一点说“缺的是意愿不是能力”，我方要在这里就把话接住。}\sub{状态说明}{，见第十一节。}
\end{evidence}
\begin{evidence}{数据}[有把握]
中国院外心脏骤停患者的存活出院率只有 1.2\%，神经功能预后良好率 0.8\%；旁观者实施心肺复苏的比例是 17.0\%，旁观者使用自动体外除颤器的比例不到 0.1\%。\sub{来源}{《中国心脏骤停与心肺复苏报告（2022年版）》，《中国循环杂志》2023年第38卷第10期1005–1017页，数据来自2020年中国七大地理区域各一个城市网点经紧急医疗服务救治的院外心脏骤停患者。}\sub{方法}{多区域城市网点登记研究。}\sub{局限}{本次备赛只打开了健康界与澎湃新闻对该报告的转述，未打开期刊原文；数字反映的是急救体系的时间瓶颈，不能直接等同于“能力缺口”。}\sub{状态说明}{。}
\end{evidence}
\begin{evidence}{案例}
切尔诺贝利，1986年4月26日凌晨。世界卫生组织记载，当时在现场的600名工作人员中，134人受到 0.8 到 16 戈瑞的极高剂量照射并出现急性放射病，其中28人在三个月内死亡；后续登记在册的善后作业人员有53万人。\sub{代表性}{这是“能力缺口”的典型而不是极端个例——不是没有人愿意去，是只能用人的身体去顶。}\sub{贴合度}{它证明的正是本论点的机制。}\sub{出处}{世界卫生组织《辐射：切尔诺贝利事故》问答页。已核实 2026-09-15。}
\end{evidence}
\closing[回扣判准]{安全那一项的净改善，而且是在没有替代方案的情况下的改善。}
\begin{clash}[对方最强攻击 → 我方回应]
攻击——现实里的瓶颈从来不是能力，是意愿和授权；卢旺达的时候大国都有能力，缺的是派不派兵。回应——我方承认意愿问题存在，这恰恰是我方论点的一部分：超级英雄的价值就在于他不需要等一个国家开会决定派不派。技术会进步，但人是在具体的那几分钟里死的，而超级英雄补的就是那几分钟。
\end{clash}
\end{fields}
\prosubsection{论点二：恶做不成}
\begin{fields}
\claim[主张]{超级英雄抬高的是“当场被拦住”的确定性，最坏的那一类恶行会因此做不成。}
\begin{mechanism}[机制]
\step 绝大多数大规模恶行的成立，依赖“不会有人当场拦住”这个预期
\step 一个有能力随时到场的存在把这个确定性整体抬高
\step 最需要靠“没人拦得住”才能实施的那一类行为（屠杀、围城、大规模强制）成本剧增。

\end{mechanism}
\begin{evidence}{学理}
威慑理论。美国国家司法研究所2016年《关于威慑的五件事》（NCJ 247350）第一条：“被抓住的确定性是远比惩罚本身更强大的威慑”；第三条明确说，把警察当作“哨兵”（sentinel）的策略特别有效，罪犯的行为更可能被“看见一个带着手铐和对讲机的警察”改变，而不是被一部提高刑罚的新法律改变。该文件内容取自 Daniel S. Nagin（2013）《二十一世纪的威慑》。一句话口头版：“让人不敢做的是会被抓住，不是抓住以后判得重。”\sub{边界}{威慑对“隐蔽的、不需要公开实施的恶”无效，我方主动承认这一点，把战场收在公开的大规模伤害上。已核实 2026-09-15。}
\end{evidence}
\begin{evidence}{数据}
Braga 等人2019年的系统综述汇总了65项研究、78次检验，把警力集中到犯罪高发的小片区域能显著降低犯罪，其中35次「总体犯罪」检验的加权平均效应量为 0.109（该综述自报的73次主效应检验总体效应量是 0.132）；而且40次关于犯罪转移的检验显示的是“犯罪控制收益向周边外溢”，不是把犯罪挤到别处去。\sub{来源}{Braga, Turchan, Papachristos \& Hureau (2019)，《实验犯罪学杂志》15卷3期289–311页，经美国国家司法研究所 CrimeSolutions 实践评级页核对。}\sub{方法}{坎贝尔协作网标准下的系统综述与元分析，27项随机实验、38项准实验，合计65项。}\sub{局限}{研究对象是有授权、有制服、可被起诉的警察，把它外推到一个不可问责的个体需要额外论证，这是反方一定会打的地方。《实验犯罪学杂志》正式版本次卡在付费墙，经同组作者同年坎贝尔综述的开放获取全文与美国国家司法研究所 CrimeSolutions 评级页核实，未打开该期刊 PDF。已核实 2026-09-18。}
\end{evidence}
\begin{evidence}{案例}
卢旺达，1994年。联合国“防止灭绝种族罪外联方案”记载，屠杀从4月6日持续到7月4日，估计有超过一百万人遇害；而联合国卢旺达援助团的兵力在4月21日从2165人减到270人。联合国自己的说法是，“国际社会的反应迟疑”使悲剧雪上加霜，联合国减轻苦难的能力“严重受制于会员国不愿加强授权、不愿增派部队”。\sub{代表性}{这是“没有人在现场挡住”的后果的典型。}\sub{出处}{联合国相关页面。已核实 2026-09-15。}
\end{evidence}
\closing[回扣判准]{安全项的净改善，而且改善的是最坏的那一档。}
\begin{clash}[对方最强攻击 → 我方回应]
攻击——历史上最大的屠杀恰恰是国家干的，而国家是最有可能收编超级英雄的主体，你的论点在最需要它的地方最不可靠。回应——正因为国家会想收编，一个不被收编的超级英雄才是屠杀最直接的障碍。对方要论证的是“他一定会被收编”，这一步对方证不了；而在他被收编之前，卢旺达那一百天里至少有人会在现场。
\end{clash}
\end{fields}
\prosubsection{论点三：人会跟着做}
\begin{fields}
\claim[主张]{超级英雄的存在不是让普通人退场，是把“这类事有人在做”变成常识，从而拉低普通人出手的门槛。}
\begin{mechanism}[机制]
\step 人做不做好事，很大程度上取决于“这件事是不是该我做、做了有没有用”
\step 一个可见的、值得效仿的榜样把这件事变成社会常识
\step 普通人自己出手的意愿上升。

\end{mechanism}
\begin{evidence}{学理}
情境启动对计划性助人行为的影响（Nelson \& Norton, 2005）。\sub{核心主张}{短暂的情境线索不仅影响当下的自发行为，也能影响人对未来助人的承诺与实际履行。}\sub{支撑}{解释了机制的第二步。}\sub{边界}{这项研究同时发现，被\textbf{具体个体}（超人）启动的人助人意愿反而下降，这一半是反方的证据，我方必须主动说出来。已核实 2026-09-15。}
\end{evidence}
\begin{evidence}{数据}
Nelson \& Norton (2005) 研究二，49名普林斯顿本科生。被要求列举“超级英雄”这个类别特征的学生，随后向一个真实校园志愿组织承诺的服务时间是每周 2.13 小时，对照组是每周 0.94 小时，两倍多。\sub{来源}{《实验社会心理学杂志》41卷4期423–430页。}\sub{方法}{随机分组的思维列举启动实验，因变量是向真实组织承诺的每周小时数。}\sub{局限}{样本只有49人，2005年的启动效应研究，在心理学可重复性危机之后需要谨慎对待；而且测的是承诺，不是最终到场。已核实 2026-09-15。}
\end{evidence}
\begin{evidence}{案例}
《中华人民共和国民法典》第一百八十四条规定，“因自愿实施紧急救助行为造成受助人损害的，救助人不承担民事责任”；第一百八十三条规定见义勇为者受损时由侵权人担责、受益人适当补偿。\sub{代表性}{这说明一个社会在制度层面认定“多一个愿意出手的人”是值得保护的公共利益，而不是麻烦。}\sub{出处}{《中华人民共和国民法典》（2020年5月28日第十三届全国人民代表大会第三次会议通过）第一百八十三条、第一百八十四条，全文见 \url{http://www.mod.gov.cn/gfbw/qwfb/yw_214049/4866121.html}。已核实 2026-09-18。}\sub{局限}{条文经中华人民共和国国防部官网转载的新华社通稿全文核实，未打开全国人大国家法律法规数据库的条文页。}
\end{evidence}
\closing[回扣判准]{自主那一项不但没被削弱，反而被抬高。}
\begin{clash}[对方最强攻击 → 我方回应]
攻击——同一篇论文的研究三直接打你：被“超人”启动的人，当场报名做志愿者的只有23\%，被“超级英雄”类别启动的有42\%；九十天后真正到场的是4\%对17\%。作者自己预测的就是“想到超人的人会比平均水平更不愿意帮忙”。回应——我方承认这个结果，而且我方认为它恰恰说明关键在于人们把他当成一个遥不可及的符号，还是当成一件有人正在做的事。研究三里的“超人”是漫画人物，不可企及；辩题假设的超级英雄是真的出现在这个世界上的人。另外请评委注意，研究三比较的是两种启动条件，作者说明省略了中性对照组，所以“4\%对17\%”不能直接读成“超级英雄让人退场”。
\end{clash}
\end{fields}
\consection{六、反方论点}
\consubsection{论点一：没人能拦}
\begin{fields}
\claim[主张]{超级英雄的出现，第一次让一种压倒性的力量落在了没有任何现存机制能够纠错的位置上。}
\begin{mechanism}[机制]
\step 能力超出所有既有约束手段
\step 什么时候出手、对谁出手，判断权由他一个人独享
\step 一旦出错、越界或被利用，没有任何现存机制能把它改回来。

\end{mechanism}
\begin{evidence}{学理}
马克斯·韦伯1919年演讲《以政治为业》。他给国家下的定义是：“一个国家是这样一个人类共同体，它在一定疆域之内（成功地）宣称对正当使用物理暴力的垄断。”紧接着的一句更关键：“使用物理暴力的权利，只在国家允许的限度之内被赋予其他机构或个人。” 一句话口头版：“现代世界的安全底座，是暴力只能经授权使用。”\sub{支撑}{说明超级英雄不是这个底座上的一个补丁，他是这个底座之外的第一个例外。}\sub{边界}{韦伯描述的是现代国家的事实特征，不等于说任何越出这条线的行为都必然有害，反方不要把它用成道德判断。已核实 2026-09-15。另配阿克顿勋爵1887年4月5日致克赖顿主教信：“权力趋于腐败，绝对权力绝对地腐败。伟人几乎总是坏人，哪怕他们行使的只是影响力而不是权威。”已核实 2026-09-15。}
\end{evidence}
\begin{evidence}{数据}
人权观察2019年报告《印度的暴力护牛》记载，2015年5月至2018年12月，印度至少44人在护牛义警袭击中被打死，其中36人是穆斯林；约280人受伤，20个邦发生100多起袭击。更关键的一条：在将近三分之一的护牛致死案件里，警方反过来对受害者或证人立案。\sub{来源}{Human Rights Watch (2019), \emph{Violent Cow Protection in India: Vigilante Groups Attack Minorities}，2019年2月发布，104页，基于2018年6月至2019年1月在四个邦的实地调查。}\sub{方法}{实地调查加案件档案，详述11起造成14人死亡的案件。}\sub{局限}{这是有强烈价值立场的国际人权组织报告；它记录的是普通人组成的义警，不是超常能力者，类比的落点必须是“能力与问责之间的落差”，不是能力大小。已核实 2026-09-15。}
\end{evidence}
\begin{evidence}{案例}
2011年10月10日凌晨，西雅图。一个真的穿上制服上街打击犯罪的人，本名本杰明·福多，自称“凤凰侠”，把一群刚从夜店出来的人当成在斗殴，朝他们喷了胡椒喷雾，随后因涉嫌袭击被捕。他说自己目击了一起未遂谋杀。西雅图警方发言人马克·贾米森说：“他打扮成什么样子都不重要，那不代表他可以得到特殊对待或者凌驾于法律之上。你不能因为你觉得别人在打架就去喷人。”最终检方因为找不到部分当事人而没有起诉。\sub{代表性}{这不是极端个例，这是“自己判断、自己动手、事后没人能算账”这条链条的完整演示。}\sub{出处}{哥伦比亚广播公司新闻2011年报道。已核实 2026-09-15。}
\end{evidence}
\closing[回扣判准]{自主项的结构性损失——世界失去的不是安全感，是出错之后改回来的能力。}
\begin{clash}[对方最强攻击 → 我方回应]
攻击——印度义警和凤凰侠都不是超级英雄，他们造成伤害恰恰是因为能力不足加判断力差，跟“超越常人的能力”没关系；而且你把“不受约束”写进定义再论证“不受约束很危险”，是循环论证。回应——我方论点不依赖“他会变坏”，我方依赖的是“我们无法验证他会不会变坏，也无法在他错的时候改回来”。类比成立的地方不是能力大小，是能力和问责之间的落差：一群没有超能力的普通人，这个落差已经足够造成四十四条人命；换成一个谁都拦不住的人，这个落差只会更大。
\end{clash}
\end{fields}
\consubsection{论点二：坏人也会有}
\begin{fields}
\claim[主张]{一旦这种力量可以落在个体身上，世界就必须按最坏的持有者来布防，而这种不对称是世界赔不起的。}
\begin{mechanism}[机制]
\step 能力的存在与持有者的品性是两件独立的事
\step 世界从此要按最坏的可能性重新布防，所有国家的第一反应是找到它、复制它、或者先下手为强
\step 对抗升级成为默认选项，而一次失手就是不可逆的。

\end{mechanism}
\begin{evidence}{学理}[有把握]
安全困境（Robert Jervis, 1978，《安全困境下的合作》，《世界政治》30卷2期167–214页）。\sub{核心主张}{一方为了自身安全而采取的手段，会降低其他人的安全，于是引发谁都不想要的对抗升级。}\sub{支撑}{解释了机制的第二步——不需要超能力真的扩散，只要它被证明可以存在，竞逐就开始了。}\sub{边界}{这个理论讲的是国家之间，用在一个个体身上需要说明“其他行为体把他当成安全威胁”这一步。}\sub{状态说明}{，见第十一节。}
\end{evidence}
\begin{evidence}{数据}
斯德哥尔摩国际和平研究所《2025年鉴》第六章记载，2025年1月全球九个拥核国家共有约12241枚核弹头，其中9614枚在可用的军事库存中，3912枚已部署在作战力量上，约2100枚处于高度戒备状态。\sub{来源}{SIPRI Yearbook 2025, ch.6 “World nuclear forces”，Hans M. Kristensen 与 Matt Korda。}\sub{方法}{基于公开资料与政府披露的年度估算。}\sub{局限}{这是类比证据，证明的不是“超能力会扩散”，而是“人类面对压倒性力量时的实际反应是什么样的”。已核实 2026-09-15。}
\end{evidence}
\begin{evidence}{案例}
1983年9月26日，莫斯科郊外的谢尔普霍夫十五号预警中心。值班的斯坦尼斯拉夫·彼得罗夫中校看到系统报告美国发射了五枚核导弹。他没有按程序上报，理由是“真要开战不会只打五枚”。后来查明，是阳光照在高空云层上的反射被卫星当成了导弹尾焰。\sub{代表性}{它说明压倒性力量的安全不取决于制度，取决于某一个人某一晚的判断。}\sub{出处}{《不列颠百科全书》“Stanislav Petrov”词条。已核实 2026-09-15。}
\end{evidence}
\closing[回扣判准]{安全项的尾部风险——一次失手足以抵消掉全部日常收益。}
\begin{clash}[对方最强攻击 → 我方回应]
攻击——辩题说的是“超级英雄的出现”，不是“超能力的普及”，你从一个人滑到所有人；而且人类造了一万两千多枚核弹，八十年里核战争发生了吗？没有，这恰恰说明人类学会了共存。回应——我方不需要超能力普及，只需要它被证明可以存在，各国的争夺就开始了，这是安全困境，不是滑坡。至于八十年没打起来，请评委记住一九八三年九月那一晚：人类和核战争之间隔着的，是一个值班军官的一次判断。“还没发生”不等于“安全”。
\end{clash}
\end{fields}
\consubsection{论点三：等靠退化}
\begin{fields}
\claim[主张]{世界会围着超级英雄重新排列，制度和普通人会一起往后退，等到他不在或者站错边的那天，世界已经没有备份。}
\begin{mechanism}[机制]
\step 有人兜底
\step 制度和个人都理性地降低投入
\step 等到需要自己站出来的那天，能力和习惯都已经不在了。

\end{mechanism}
\field{学理与数据（同一来源，本论点的核心）}{Nelson \& Norton (2005) 研究三，112名普林斯顿本科生，随机分成两组，一组列举“超人”的特征，一组列举“超级英雄”这个类别的特征。结果：被“超人”启动的人，当场报名做志愿者的只有23\%（57人中13人），被“超级英雄”类别启动的有42\%（55人中23人）。九十天后真的到场参加组织会议的，前者是4\%（56人中2人），后者是17\%（52人中9人），差了四倍。作者事先的预测就是“被超人启动的人会比平均水平更不愿意帮忙”。\sub{来源}{《实验社会心理学杂志》41卷4期423–430页。}\sub{方法}{随机分组启动实验，因变量是真实的报名率与九十天后的实际到场率。}\sub{局限}{作者说明研究三省略了中性对照组，理由是前几项研究显示对照组落在两者之间，所以这两个数字之间的差距不能直接读成“相对于没有超级英雄的世界”；样本是2005年的普林斯顿本科生；启动效应研究在心理学可重复性危机之后需要谨慎。正方一定会打这一点，我方要主动说明并把制度侧证据提到前面。已核实 2026-09-15。}
}
\begin{evidence}{数据}
瑞典哥德堡大学 V-Dem 研究所《2025年民主报告》：全球有45个国家正处在威权化进程中；这45个国家里有27个在进程开始时还是民主国家，到2024年只剩9个仍然是民主国家，报告把这个比例称为67\%的“死亡率”。另外，全球约72\%的人口生活在威权国家，是1978年以来最高。\sub{来源}{V-Dem Institute (2025), \emph{Democracy Report 2025: 25 Years of Autocratization – Democracy Trumped?}。}\sub{方法}{基于全球超过四千名国别专家编码的 V-Dem 数据集与“民主侵蚀与韧性”数据集。}\sub{局限}{这是关于政治制度的证据，用在超级英雄身上是类比；它证明的是“一旦一个社会开始依赖某个不受制衡的角色，制度往往回不来”，不是“超级英雄会当独裁者”。已核实 2026-09-15。}
\end{evidence}
\begin{evidence}{案例}
守护天使队，1979年由柯蒂斯·斯利瓦在纽约创立，成员戴红色贝雷帽在地铁巡逻。据斯利瓦自己的说法，一年之内成员从13人扩到1000人。1992年，斯利瓦承认该组织六起“成功打击犯罪”的事迹是编造的，包括所谓营救劫案和强奸受害者；同年纽约分部只剩下大约30到125名成员。\sub{代表性}{它展示的是“被依赖的英雄叙事”如何先膨胀再塌陷，而这中间真实的治安改善很难被验证。}\sub{出处}{《不列颠百科全书》“Guardian Angels”词条。已核实 2026-09-15。}
\end{evidence}
\closing[回扣判准]{自主项的直接损失，以及长期安全项的隐性损失。}
\begin{clash}[对方最强攻击 → 我方回应]
攻击——制度退化的前提是他可靠到让人敢依赖，可你的论点一又说他不可靠，你自相矛盾。回应——不矛盾。人依赖的是“感觉上有人兜底”，不是“经过验证的可靠”。V-Dem 里那27个国家的选民，也不是先验证了强人可靠才把制度交出去的。人先松手，再发现没人接着，这个顺序才是危险所在。
\end{clash}
\end{fields}
\section{七、正方反驳库（针对反方全部可能论点）}
\begin{rebuttals}
\rebuttal{1}{没人能拦：一种不可问责的力量出现，世界失去纠错能力}{你说的纠错能力，在他出现之前对那些事本来也不存在}{打机制：卢旺达那一百天，现存制度既没有拦住屠杀，也没有事后把一百万人改回来。所谓“纠错能力”在最需要它的场合本来就是空的。打判准：判准问的是净变化，不是他是否完美。对方要证明的是“有他之后比没他更糟”，不是“他不如一个理想制度”}{追问：那你承认他不可问责吗？→ 承认，我方定义里就写着。但不可问责不等于灾难，消防员冲进火场时也没有人能当场问责他的每一个判断，我们接受它是因为净收益为正}{最终}
\rebuttal{2}{坏人也会有：风险不对称，一次失手抵消全部收益}{辩题问的是超级英雄的出现，不是超能力的普及，你这一步没有论证}{打机制：从“一个人有”推到“坏人也会有”，中间缺了能力如何复制这一步，对方拿不出任何依据。打论据：核弹是国家用工业体系造出来的技术，可以被条约计数和约束；一个人不是技术。而且八十年过去，一万两千多枚核弹没有导致核战争，这个数据恰恰说明人类有能力与压倒性力量共存}{追问：那一九八三年九月那一晚呢？→ 那一晚说明的是自动化预警系统会误判，而拦住它的恰恰是一个人的判断。对方的例子里，救了世界的是人，不是制度}{最终}
\rebuttal{3}{等靠退化：制度和普通人一起往后退}{你引用的那项研究里，作者自己说明没有中性对照组}{打论据：研究三比较的是“想具体的超人”和“想超级英雄这个类别”，不是“有”和“没有”，所以四倍差距不能读成“超级英雄让人退场”。而同一篇论文的研究二里，被“超级英雄”类别启动的人承诺的志愿时间是对照组的两倍多。打机制：一个真实出现在世界上的人，会让“这类事有人在做”成为常识，这和一个漫画符号不是一回事}{追问：那你承认想到超人的人更不愿意帮忙吗？→ 我方承认这个结果。我方认为它说明的是符号和真人的区别，而这正是辩题的前提：他真的出现了}{最终}
\rebuttal{4}{超级英雄会被国家收编，变成武器}{对方要论证的是“一定会”，这一步对方证不了}{打举证责任：这是一个关于未来的强断言，对方没有任何依据说明收编必然成功。打机制：真正无法拒绝收编的是技术和武器，而一个有自己判断的人恰恰是最难被收编的那一类。反过来说，正因为国家想收编他，他不被收编时才是屠杀最直接的障碍}{追问：曼哈顿计划呢？→ 那是一项技术，不是一个会说不的人}{常见}
\rebuttal{5}{他的出现会招来更强的反派，净安全下降}{这是用设定当论据，不是用现实当论据}{打论据：反派因英雄而生是漫画的叙事惯例，是编剧为了持续连载设计的，不是关于现实的命题。请对方给出一条现实世界的证据。打机制：现实里的威慑研究结论正相反，被抓住的确定性上升，公开实施的恶行会减少}{追问：军备竞赛不就是这样吗？→ 军备竞赛发生在对称的行为体之间，而本题假设的是一个非对称的个体}{常见}
\rebuttal{6}{他一个人的判断代替了所有人的商量，取消了民主}{他做的是紧急救援，不是立法}{打定义：我方定义里他介入的是公共安全事务，不是公共决策。打机制：救一个正在被打的人，不需要先开会表决；对方把“当场制止”说成“替代治理”}{追问：他决定去救谁，这本身就是分配 → 那是选择，不是权力。任何一个先到现场的医生都在做这种选择}{淘汰}
\rebuttal{7}{人会崇拜他，形成个人崇拜和造神}{崇拜的对象是他做的事，不是他的位置}{打机制：造神的条件是他垄断了资源分配与信息，而超级英雄并不掌握这两样。打论据：对方用的政治学证据讲的是掌权者，不是救援者}{追问：V-Dem 那个67\%呢？→ 那27个国家交出去的是选票和法院，不是被人从火里背出来}{常见}
\rebuttal{8}{他改变不了结构性问题，贫困、不平等、气候，他只能打人}{判准问的是净变化，不是他能不能解决所有问题}{打判准：没有任何一项改善需要证明它能解决一切。打机制：他补的是最急、最没有替代方案的那一段，这一段恰恰是制度最难覆盖的}{追问：那他的作用很有限 → 有限不等于负面，判准比的是有和没有}{常见}
\rebuttal{9}{受害者只能等待被拯救，人的主体性被剥夺}{被救下来的人，才有机会谈主体性}{打机制：主体性的前提是活着。打论据：《民法典》第一百八十四条把见义勇为写进法律，说明社会认定多一个愿意出手的人是好事，而不是对旁人主体性的剥夺}{追问：那为什么要立这一条？→ 因为社会担心的是没人敢出手，不是有人出手}{淘汰}
\rebuttal{10}{他会死、会老、会被打败，世界赌在一个人身上}{世界从来不是靠他活着，是多了一层他}{打机制：他不在的时候世界回到原来的样子，那是零，不是负数。对方要证明的是低于零。打判准：净变化在他出现的那些年里是正的，这些年里被救下来的人不会因为他后来不在了就重新死一次}{追问：那退化呢？→ 见第3条}{常见}
\end{rebuttals}
\section{八、反方反驳库（针对正方全部可能论点）}
\begin{rebuttals}
\rebuttal{1}{救得下来：他补上了现有制度补不上的能力缺口}{现实里的瓶颈从来不是能力，是授权和意愿}{打机制：卢旺达的时候，能力是现成的，大国的军队随时可以到，缺的是派不派兵这个决定；联合国自己的说法是受制于会员国不愿加强授权、不愿增派部队。切尔诺贝利也不是没人愿意去，是当时的机械在高辐射下失灵，这是技术问题，技术在进步。打判准：把这些换成一个人，缺口没有被补上，只是从“没人决定”变成“一个人决定”}{追问：那技术进步之前死的人怎么算？→ 我方从不否认有人救不下来，我方的问题是：把决定权交给一个谁都管不了的人，是不是一个更好的办法}{最终}
\rebuttal{2}{恶做不成：他抬高了被当场拦住的确定性}{你引用的威慑研究，讲的全是受授权、可被起诉的警察}{打论据：国家司法研究所说的“哨兵”是穿制服、带编号、能被投诉和起诉的人，热点警务的65项研究研究的也是这种人。把它外推到一个不可问责的个体，你把“被抓住的确定性”偷换成了“某个人心情的确定性”。打机制：同一套威慑逻辑还会把对手推向更极端、更隐蔽的手段}{追问：难道威慑不起作用吗？→ 起作用，正因为它起作用，作用的方向由他一个人决定，这才是问题}{最终}
\rebuttal{3}{人会跟着做：榜样把普通人出手的门槛拉低}{你引用的那篇论文，另一半直接写着相反的结果}{打论据：同一篇论文的研究三，112个人，被“超人”启动的人当场报名志愿只有23\%，被“超级英雄”类别启动的有42\%；九十天后真正到场的是4\%对17\%，差了四倍。作者事先预测的就是这个方向。你只引了研究二。打机制：人和一个不可企及的榜样比较之后，会觉得自己做的那点事没有意义}{追问：研究三没有对照组 → 是的，我方承认。所以我方把它当作机制的说明，主证据是制度侧的数据：45个威权化国家里有27个起点是民主国家，只剩9个还是}{最终}
\rebuttal{4}{他能解决人类靠自己解决不了的问题，比如小行星、瘟疫}{这是设定，不是论据}{打论据：请对方给出一条现实世界的依据说明这类能力存在。打判准：即使存在，判准同时要问代价，而这类能力越大，不可问责的问题越严重}{追问：那你承认能力是好事吗？→ 能力本身中性，落在没人能管的地方才是问题}{常见}
\rebuttal{5}{有他兜底，普通人反而更敢站出来}{数据的方向正相反}{打论据：见第3条的四倍差距。打机制：兜底改变的是“该不该我做”的判断，人退后一步是理性的}{追问：那消防员在场的时候人就不救人了吗？→ 消防员是可以被呼叫、被问责的公共服务，他是不是来、来不来得及是可预期的，这和一个凭自己判断出现的人不一样}{淘汰}
\rebuttal{6}{世界需要一个明确的善的样子，他提供了道德坐标}{道德坐标交给一个不可问责的人，就是把善的定义也交出去了}{打机制：坐标意味着标准，标准由他一个人给，这恰恰是我方论点一说的事。打论据：守护天使队1992年承认六起“英雄事迹”是编造的，被崇拜的叙事本身就会被制造}{追问：那英雄故事没有价值吗？→ 有价值，前提是讲故事的人和被讲的人不是同一个}{淘汰}
\rebuttal{7}{他的存在会倒逼人类科技与制度进步}{倒逼的方向是竞逐，不是进步}{打学理：安全困境说的正是这个，一方为自身安全采取的手段会降低其他人的安全。打论据：人类上一次面对压倒性力量，倒逼出来的是一万两千多枚核弹头}{追问：核威慑不是维持了和平吗？→ 一九八三年九月那一晚，维持和平的是一个军官不按程序办事}{常见}
\rebuttal{8}{风险可以靠契约、登记、监督来解决}{谁去执行}{打机制：所有监督手段的最终依托都是强制力，而他恰恰是强制力管不了的那个人。签了字不遵守的时候，下一步是什么？对方答不出这一步。打学理：韦伯说得很清楚，使用物理暴力的权利只在国家允许的限度内被赋予个人}{追问：那舆论和道德呢？→ 舆论对一个不需要任何人配合的人没有约束力}{常见}
\rebuttal{9}{灾难必须是已经发生的重大恶化，你只证明了风险}{我方证明的不是风险，是已经发生的结构改变}{打判准：从他出现那天起，世界就进入了“最终手段失效”的状态，这不是概率，是已经成立的事实。打机制：这就像一个国家的法院从此对某个人无效，不需要等他犯罪，法治已经缺了一块}{追问：那什么都没发生呢？→ 纠错能力的消失本身就是后果}{常见}
\rebuttal{10}{现实里的义警和他不一样，你的类比不成立}{类比的落点是能力和问责之间的落差，不是能力大小}{打论据：印度的护牛义警没有任何超常能力，这个落差已经造成至少四十四条人命，将近三分之一的案子里警方还反过来追究受害者。换成一个谁都拦不住的人，落差只会更大。打机制：对方其实在替我方论证——能力越强，落差越大}{追问：可他是好人 → 我方从不假设他是坏人。我方的问题是我们无法验证，也无法在他错的时候改回来}{常见}
\end{rebuttals}
\section{九、持方优劣评估}
\begin{assessment}
\assess{定义空间}{常识里的超级英雄是“救人的好人”，表面看偏向正方。但常识里同样包含“他不归警察管”这一条，而这一条一旦进入中立定义，反方的整个战场就自动打开。正方若试图把它删掉，会被指定义杀；若保留，就要一直防守。反方只要守住这一条，全场都有话说}{偏反}
\assess{论证义务}{“更是”要求双方都证明主导性，两边义务对称。但正方的收益是可数的（被救的人），反方的代价是结构性的、不可数的（纠错能力），评委更容易给可数的东西记分}{偏正}
\assess{论据可得性}{现实世界里“不受约束的执法者”的经验证据非常丰富：义警研究、私刑统计、个人化统治数据，都是现成的。而“超常能力带来净收益”几乎没有直接对应物，正方只能用反面案例（切尔诺贝利、卢旺达），而这些案例反方可以反手说“缺的是授权不是能力”}{偏反}
\assess{评委直觉}{新生赛评委的第一反应多半是“有超人当然好啊”。正方起手有直觉红利，反方必须在第一分钟内让评委意识到“那谁来管他”}{偏正}
\assess{赛制顺序}{正方先立论，定义主动权在正方手里，这是正方的结构优势。但反四先结辩可以封路，而这道题的最后一块最可能落在价值层，封路的价值高。两项大致抵消}{均衡}
\end{assessment}
\begin{fields}
\claim[结论]{\textbf{正方（更是幸福）较劣势}，劣势主要来自定义空间与论据可得性这两个维度。正方的直觉红利只值前一分钟；一旦反方把“谁来监督”抛出来，正方就会被迫全场防守。}
\begin{procard}{劣势方（正方）打法}\end{procard}
\begin{fields}
\field{抢先定义}{正一必须自己把“不受既有授权与问责”说出来。藏着不说，反四第一个问题就会把它挖出来，那时候正方是“被抓到”，自己说是“坦荡”。说完立刻用判准框住：不受授权不等于灾难，判准问的是净变化。
}
\field{收窄战场}{只打“最坏的那一类伤害”和“具体的时间窗口”，不接“日常治安该不该由超人管”“他会不会当总统”这些战场，一句话带回主战场。
}
\field{判准之争}{全场反复问反方同一个问题——\textbf{你说的灾难，是和哪个世界比？} 如果是和一个能及时救人的完美制度比，那个世界不存在。这个问题要在正一立论里就抛出来，质询、盘问、自由辩各追一次。
}
\field{价值层}{正四结辩的升华落在“能力不到的人在现场的那几分钟”，不要谈文明与时代。
}
\end{fields}
\begin{concard}{优势方（反方）要防的}\end{concard}
\begin{enumerate}[leftmargin=1.8em,itemsep=2pt,topsep=2pt]
\item 最容易翻车的地方是滑向“超级英雄一定会变坏”。反方的论点不依赖他变坏，只依赖“我们无法验证、也无法改回来”。一旦滑过去，正方一句“你凭什么假设他是坏人”就能打回来。
\item 第二个坑是把漫画情节当论据。评委会立刻扣分，而且正方一说“这是设定”，反方的整套证据链就显得不严肃。
\item 第三个坑是引用 Nelson \& Norton 研究三时忘了说没有对照组。正方一旦抓到这一点并且指出反方回避了局限，反方的诚信分会掉。
\end{enumerate}
\end{fields}
\section{十、口径表}
\prosubsection{10.1 正方口径表}
\begin{canon}{}
\canongroup{定义}
\canonrow{超级英雄}{拥有远超常人的能力，主动用这种能力介入公共安全事务，行动依据是自己的判断，不由现有法律和权力机构授权，也不由它们问责}{“能力超常、主动救人、不受现有权力管辖的人”}{“定义上永远做对事的人”；“就是特别勇敢的普通人”}{本文档 3.1；《不列颠百科词典》“superhero”词条}
\canonrow{出现}{这种人存在于世界上成为既定事实，并且被世界知道}{“他在这个世界上，而且大家都知道”}{“他在现场出手救人”（把出现缩成一次救援）}{本文档 3.2}
\canonrow{灾难}{世界整体处境出现重大且难以逆转的恶化}{“整个世界变得更糟，而且回不去”}{“只要有风险就是灾难”}{本文档 3.3}
\canongroup{判准}
\canonrow{中立判准}{有超级英雄的世界和没有超级英雄的世界放在一起，长期看普通人的处境是变好还是变坏；处境包含安全和自主两项}{“净变化落在哪一边”}{“只看救了多少人”}{本文档第四节}
\canongroup{论点}
\canonrow{救得下来}{有一类伤害的瓶颈是能力本身，超级英雄补的正是这个缺口}{“有些事不是想不想，是能不能”}{“他能解决一切问题”}{本文档 5.1}
\canonrow{恶做不成}{他抬高了当场被拦住的确定性，最坏的那一类恶行会因此做不成}{“让最坏的事做不成”}{“犯罪会消失”；“他能预防所有犯罪”}{本文档 5.2}
\canonrow{人会跟着做}{他把“这类事有人在做”变成常识，普通人出手的门槛因此降低}{“多一个人做，就多一些人敢做”}{“有他在大家就都会变成好人”}{本文档 5.3}
\canongroup{数据}
\canonrow{切尔诺贝利}{现场600人中134人出现急性放射病，28人在三个月内死亡，善后作业人员53万人}{“一百多人当场受到极高剂量照射，二十八个人三个月内没了”}{把28人说成“上百人”；把53万说成“上百万”}{世界卫生组织《辐射：切尔诺贝利事故》问答页，已核实 2026-09-15}
\canonrow{卢旺达}{1994年4月6日至7月4日，估计超过一百万人遇害；联合国卢旺达援助团兵力4月21日从2165人减到270人}{“一百天，一百多万人；维和部队从两千多人减到不到三百人”}{说成“联合国见死不救”（联合国原文是会员国不愿加强授权与增派部队）}{联合国防止灭绝种族罪外联方案历史背景页，已核实 2026-09-15}
\canonrow{威慑确定性}{美国国家司法研究所2016年《关于威慑的五件事》：被抓住的确定性是远比惩罚本身更强大的威慑}{“让人不敢做的是会被抓住，不是判得重”}{说成“警察越多犯罪越少”}{NIJ (2016), NCJ 247350，依据 Nagin (2013)，已核实 2026-09-15}
\canonrow{热点警务}{Braga 等2019年系统综述，65项研究78次检验，总体犯罪加权平均效应量0.109，且是收益外溢而非犯罪转移}{“六十多项研究汇总下来，犯罪确实降了，周边地区也跟着降”}{说成“降低了16\%”或任何未经核实的百分比}{Braga et al. (2019)，《实验犯罪学杂志》15(3):289–311，经 NIJ CrimeSolutions 核对，已核实 2026-09-15}
\canonrow{榜样效应}{Nelson \& Norton 2005 研究二，49人，被“超级英雄”类别启动的学生承诺每周2.13小时志愿服务，对照组0.94小时}{“承诺的志愿时间是对照组的两倍多”}{隐瞒研究三的相反结果}{《实验社会心理学杂志》41(4):423–430，已核实 2026-09-15}
\canongroup{案例}
\canonrow{好人条款}{《民法典》第一百八十四条：因自愿实施紧急救助行为造成受助人损害的，救助人不承担民事责任}{“法律专门给见义勇为的人免责”}{说成“法律鼓励所有人去冒险”}{《中华人民共和国民法典》（2020），已核实 2026-09-15}
\canongroup{承认}
\canonrow{主动坦白项}{我方承认超级英雄不受现有权力问责；我方承认 Nelson \& Norton 研究三的结果与我方论点三方向相反}{“这一点我方承认”}{回避、否认、说对方引错了}{本文档 5.3 与第七节第3条}
\end{canon}
\consubsection{10.2 反方口径表}
\begin{canon}{}
\canongroup{定义}
\canonrow{超级英雄}{与正方同一版中立定义：能力超常、主动介入、不受既有授权与问责}{“谁都拦不住、谁都管不了的那个人”}{“他就是个不受约束的暴力持有者”（删掉英雄二字）}{本文档 3.1}
\canonrow{出现}{这种人存在于世界上成为既定事实，并且被世界知道}{“从他出现那天起，世界就知道有这么个人”}{“超能力会普及到所有人身上”}{本文档 3.2}
\canonrow{灾难}{世界整体处境出现重大且难以逆转的恶化，核心是纠错能力的丧失}{“错了以后改不回来”}{“只要有风险就是灾难”}{本文档 3.3}
\canongroup{判准}
\canonrow{中立判准}{同正方，安全加自主两项的长期净变化}{“长期看普通人的处境”}{“只看最坏情况”}{本文档第四节}
\canongroup{论点}
\canonrow{没人能拦}{一种压倒性力量落在了没有任何现存机制能够纠错的位置上}{“出事以后没人能把它改回来”}{“他迟早会作恶”；“绝对权力必然腐败所以他必然腐败”}{本文档 6.1}
\canonrow{坏人也会有}{世界从此必须按最坏的持有者来布防，而这种不对称赔不起}{“只要它能存在，世界就得按最坏的来准备”}{“超能力一定会扩散”}{本文档 6.2}
\canonrow{等靠退化}{制度和普通人会一起往后退，等他不在时世界已经没有备份}{“人先松手，再发现没人接着”}{“有他在人就全都躺平了”}{本文档 6.3}
\canongroup{数据}
\canonrow{印度护牛义警}{人权观察2019年报告：2015年5月至2018年12月至少44人被打死，其中36人是穆斯林，约280人受伤，20个邦100多起袭击；近三分之一的致死案件里警方反过来对受害者或证人立案}{“四十四条人命，三分之一的案子警察反过来查受害者”}{把44说成“上百人”；说成“印度政府组织的屠杀”}{Human Rights Watch (2019)，已核实 2026-09-15}
\canonrow{美国私刑}{平等正义倡议组织记录：1877年至1950年，南方十二州有4084起种族恐怖私刑，其他州300多起，比此前最全面的统计多出至少800起}{“四千多起私刑，比以前统计的还多出八百多起”}{把它说成“现在还在发生”}{Equal Justice Initiative, \emph{Lynching in America}，已核实 2026-09-15}
\canonrow{核弹头}{斯德哥尔摩国际和平研究所2025年鉴：2025年1月全球约12241枚核弹头，9614枚在可用军事库存中，3912枚已部署，约2100枚处于高度戒备}{“一万两千多枚核弹头”}{说成“超能力会像核武器一样扩散”}{SIPRI Yearbook 2025, ch.6，已核实 2026-09-15}
\canonrow{制度塌陷}{V-Dem《2025年民主报告》：45个国家正在威权化，其中27个起点是民主国家，到2024年只剩9个仍是民主国家；全球约72\%人口生活在威权国家}{“二十七个国家交出去，只有九个还回得来”}{说成“超级英雄会当独裁者”}{V-Dem Institute (2025)，已核实 2026-09-15}
\canonrow{退化实验}{Nelson \& Norton 2005 研究三，112人：被“超人”启动的当场报名率23\%，被“超级英雄”类别启动的42\%；九十天后实际到场率4\%对17\%}{“想到超人的那一组，三个月后到场的人少了四倍”}{隐瞒“研究三没有中性对照组”这一局限}{《实验社会心理学杂志》41(4):423–430，已核实 2026-09-15}
\canongroup{案例}
\canonrow{凤凰侠}{2011年10月西雅图，本名本杰明·福多的“凤凰侠”把一群刚出夜店的人当成斗殴，喷了胡椒喷雾后被捕；警方说打扮成什么样都不能凌驾于法律之上；检方最终未起诉}{“一个真的穿制服上街打击犯罪的人，喷错了人”}{说成“所有想做好事的人都会伤人”}{哥伦比亚广播公司新闻2011年报道，已核实 2026-09-15}
\canonrow{守护天使队}{1979年纽约创立，一年内成员从13人到1000人；1992年创始人承认六起打击犯罪事迹是编造的，同年纽约分部只剩约30到125人}{“六起英雄事迹是编的”}{说成“所有民间巡逻都是骗局”}{《不列颠百科全书》“Guardian Angels”词条，已核实 2026-09-15}
\canongroup{学理}
\canonrow{暴力垄断}{韦伯1919年《以政治为业》：国家是在一定疆域内成功宣称对正当使用物理暴力之垄断的人类共同体；使用物理暴力的权利只在国家允许的限度内被赋予其他机构或个人}{“现代世界的安全底座，是暴力只能经授权使用”}{把它当成道德判断（“越出这条线就是恶”）}{韦伯《以政治为业》（1919），已核实 2026-09-15}
\canongroup{承认}
\canonrow{主动坦白项}{我方承认没有超级英雄的世界里确实有人救不下来；我方承认 Nelson \& Norton 研究三没有中性对照组}{“这一点我方承认”}{否认救援价值；隐瞒研究局限}{本文档 6.3}
\end{canon}
\section{十一、论据出处清单}
\begin{sourcelist}
\source{1}{超级英雄词典义两条}{Britannica Dictionary，“superhero” 词条，\url{https://www.britannica.com/dictionary/superhero}，2026 年 9 月 18 日访问}{a fictional character who has amazing powers (such as the ability to fly)……a very heroic person}{2026-09-18}{已核实}{双方定义}{}
\source{2}{研究二：2.13 vs 0.94 小时/周，N=49}{Nelson, L. D., \& Norton, M. I. (2005). \emph{From student to superhero: Situational primes shape future helping}. \emph{Journal of Experimental Social Psychology} 41(4): 423–430，DOI 10.1016/j.jesp.2004.08.003；哈佛商学院公开全文 PDF \url{https://www.hbs.edu/ris/download.aspx?name=nelson+norton+from+student+to+superhero.pdf}}{Forty-nine Princeton undergraduates……participants who had been primed with the category superhero volunteered more than twice as many hours (M = 2.13 h/week) as participants in the control condition (M = .94 h/week)}{2026-09-18}{已核实}{正方论点三}{原文见第 425–426 页 Study 2 的 Method 与 Results and discussion 两段}
\source{3}{研究三：23\% vs 42\%；九十天后 4\% vs 17\%，N=112}{Nelson, L. D., \& Norton, M. I. (2005). \emph{From student to superhero: Situational primes shape future helping}. \emph{Journal of Experimental Social Psychology} 41(4): 423–430，DOI 10.1016/j.jesp.2004.08.003；哈佛商学院公开全文 PDF \url{https://www.hbs.edu/ris/download.aspx?name=nelson+norton+from+student+to+superhero.pdf}}{people were more likely to volunteer when primed with superhero (42\%, 23 of 55) than when primed with Superman (23\%, 13 of 57)……primed with superhero were more likely to show up to participate in the group (17\%, 9 of 52) than were people that had been primed with Superman (4\%, 2 of 56)}{2026-09-18}{已核实}{反方论点三}{九十天到场率那一检验原文写的是 $\chi^2$ (1, N = 108)，与研究三总样本 112 人不同，是因为有 4 人未纳入该项分析；「研究三省略中性对照组」见同文 Study 3 讨论段}
\source{4}{印度护牛义警：44死、36名穆斯林、约280伤、20邦100余起、近三分之一案件警方反查受害方}{Human Rights Watch (2019), \emph{Violent Cow Protection in India: Vigilante Groups Attack Minorities}，2019 年 2 月 18 日发布，\url{https://www.hrw.org/report/2019/02/18/violent-cow-protection-india/vigilante-groups-attack-minorities}}{Between May 2015 and December 2018, at least 44 people—36 of them Muslims—were killed across 12 Indian states……around 280 people were injured in over 100 different incidents across 20 states}{2026-09-18}{已核实}{反方论点一}{警方反查受害方一条见同页 In almost a third of the cow-related vigilante killings since 2015, police filed cases against victims or witnesses；注意 44 人遇害发生在 12 个邦，20 邦 100 余起对应的是受伤与袭击起数，两个口径不要混说}
\source{5}{美国种族恐怖私刑 4084 起（南方十二州）+ 300 余起（其他州），1877–1950}{Equal Justice Initiative, \emph{Lynching in America: Confronting the Legacy of Racial Terror}（第三版，2017 年增补，2020 年 9 月 16 日版 PDF），\url{https://eji.org/wp-content/uploads/2020/09/lynching-in-america-3d-ed-091620.pdf}}{EJI has documented 4084 racial terror lynchings in twelve Southern states between the end of Reconstruction in 1877 and 1950……also documented more than 300 racial terror lynchings in other states during this time period}{2026-09-18}{已核实}{反方论点一备用}{报告正文另一处把「其他州」的数字写成 341 起，集中在 Illinois、Indiana、Kansas、Maryland、Missouri、Ohio、Oklahoma、West Virginia 八州；赛场只说「300 余起」}
\source{6}{45国威权化；27个起点为民主、仅9个仍是民主；72\%人口在威权国家}{V-Dem Institute (2025), \emph{Democracy Report 2025: 25 Years of Autocratization – Democracy Trumped?}，\url{https://www.v-dem.net/documents/60/V-dem-dr__2025_lowres.pdf}}{27 of the 45 autocratizers were democracies at the start of their episode. Only 9 remain democracies in 2024. The fatality rate is 67\%.……Nearly 3 out of 4 persons in the world – 72\% – now live in autocracies}{2026-09-18}{已核实}{反方论点三}{67\% 报告原文用的词是 fatality rate（不是 mortality rate）；同段另写明 27 个中已有 18 个变成威权国家}
\source{7}{12241枚核弹头（2025年1月）、9614枚可用、3912枚部署、约2100枚高度戒备}{SIPRI Yearbook 2025, ch.6 “World nuclear forces”（Hans M. Kristensen \& Matt Korda），章节 PDF 经 SIPRI 官网 \url{https://www.sipri.org/yearbook/2025/06} 下载核对}{Of the total global inventory of an estimated 12 241 warheads in January 2025, about 9614 were in military stockpiles and available for potential use……An estimated 3912 of those warheads were deployed with missiles and aircraft……Approximately 2100 of the deployed warheads were kept in a state of high operational alert on ballistic missiles}{2026-09-18}{已核实}{反方论点二}{}
\source{8}{被抓住的确定性远强于惩罚严厉性；警察作为“哨兵”特别有效}{National Institute of Justice (2016), \emph{Five Things About Deterrence}, NCJ 247350，2016 年 5 月，\url{https://www.ojp.gov/pdffiles1/nij/247350.pdf}；依据 Nagin (2013), \emph{Crime and Justice} vol.42}{The certainty of being caught is a vastly more powerful deterrent than the punishment……Strategies that use the police as ‘sentinels,’ such as hot spots policing, are particularly effective}{2026-09-18}{已核实}{正方论点二}{同页紧接一句可一并背下：A criminal's behavior is more likely to be influenced by seeing a police officer with handcuffs and a radio than by a new law increasing penalties}
\source{9}{热点警务：65项研究78次检验，其中35次总体犯罪检验效应量0.109，收益外溢非犯罪转移}{Braga, Turchan, Papachristos \& Hureau (2019), \emph{Journal of Experimental Criminology} 15(3):289–311（Springer 页面本次 403，付费墙）；同组作者同年坎贝尔系统综述全文开放获取 \emph{Campbell Systematic Reviews} 15(3):e1046，DOI 10.1002/cl2.1046，\url{https://pmc.ncbi.nlm.nih.gov/articles/PMC8356500/}；0.109 一数经美国国家司法研究所 CrimeSolutions 实践评级页 \url{https://crimesolutions.ojp.gov/ratedpractices/hot-spots-policing} 核对}{Overall, aggregating the results from 35 independent tests, Braga and colleagues (2019) found an overall statistically significant weighted mean effect size of 0.109 in favor of hot spots policing}{2026-09-18}{已核实}{正方论点二}{65 项与 78 次的原句是坎贝尔综述的 Sixty-five studies containing 78 tests of hot spots policing interventions were identified；该综述自报的 73 次主效应检验总体效应量是 0.132，0.109 是 CrimeSolutions 就「总体犯罪」这一结果类别汇总 35 次独立检验得出的，两个数不要混说；40 次转移／外溢检验的平均效应量为 0.086，方向是收益外溢}
\source{10}{卢旺达：逾百万人遇害；UNAMIR 由2165人减至270人；联合国自述会员国不愿增授权增兵}{联合国「1994 年卢旺达对图西人灭绝种族罪外联方案」历史背景页，\url{https://www.un.org/en/preventgenocide/rwanda/historical-background.shtml}，2026 年 9 月 18 日访问}{more than one million people are estimated to have perished……On 21 April, after other countries asked to withdraw troops, the UNAMIR force reduced from an initial 2,165 to 270}{2026-09-18}{已核实}{正方论点二}{会员国不愿增授权增兵一条是同页的表述，赛场只说「联合国自己的说法是会员国不愿意扩大授权、不愿意增兵」，不要报成引号原话}
\source{11}{切尔诺贝利：现场600人中134人急性放射病、28人三个月内死亡、53万善后作业人员}{世界卫生组织《Radiation: The Chernobyl accident》问答页，\url{https://www.who.int/news-room/questions-and-answers/item/radiation-the-chernobyl-accident}，2026 年 9 月 18 日访问}{Of 600 workers present on the site during the early morning of 26 April 1986, 134 received very high doses (0.8-16 Grey) and suffered from acute radiation sickness. Of those, 28 workers died in the first three months}{2026-09-18}{已核实}{正方论点一}{53 万一数见同页 Among 530,000 registered recovery operation workers who worked between 1986 and 1990, the average dose was 120 mSv}
\source{12}{1983年9月26日苏联预警误报五枚导弹，起因是阳光在高云上的反射}{Encyclopaedia Britannica, “Stanislav Petrov” 词条，\url{https://www.britannica.com/biography/Stanislav-Petrov}，2026 年 9 月 18 日访问}{Five nuclear missiles in total, straight from an American base in North Dakota……The technology, called Oko satellites, mistakenly perceived sunlight on high altitude clouds as ballistic missile engine exhaust}{2026-09-18}{已核实}{反方论点二}{}
\source{13}{守护天使队：1979年创立、一年内13人到1000人、1992年承认六起事迹编造、同年纽约分部约30至125人}{Encyclopaedia Britannica, “Guardian Angels” 词条，\url{https://www.britannica.com/topic/Guardian-Angels}，2026 年 9 月 18 日访问}{from 13 to 1,000…within a period of a year}{2026-09-18}{已核实}{反方论点三}{1992 年承认六起打击犯罪事迹系编造、以及同年纽约分部成员约 30 至 125 人，见同词条讲 1992 年的那一段；创始人为 Curtis Sliwa，1979 年创立于纽约布朗克斯}
\source{14}{凤凰侠（本杰明·福多）2011年10月胡椒喷雾事件与西雅图警方表态}{CBS News，\emph{Seattle “superhero” Phoenix Jones arrested following pepper-spray attack}，2011 年 10 月 14 日，\url{https://www.cbsnews.com/news/seattle-superhero-phoenix-jones-arrested-following-pepper-spray-attack/}}{Just because he's dressed up in costume, it doesn't mean he's in special consideration or above the law. You can't go around pepper spraying people because you think they are fighting.}{2026-09-18}{已核实}{反方论点一}{警方表态人为西雅图警局发言人 Det. Mark Jamieson；同报道写明福多称对方在斗殴、警方称对方是在跳舞}
\source{15}{国家对正当暴力的垄断；个人使用暴力的权利只在国家允许限度内}{Max Weber, \emph{Politics as a Vocation}（1919 年演讲），牛津大学贝利奥尔学院公开节选 PDF \url{https://www.balliol.ox.ac.uk/sites/default/files/politics_as_a_vocation_extract.pdf}}{a state is a human community that (successfully) claims the monopoly of the legitimate use of physical force within a given territory……the right to use physical force is ascribed to other institutions or to individuals only to the extent to which the state permits it}{2026-09-18}{已核实}{反方论点一}{核实用的是通行英译节选（Gerth 与 Mills 译本），不是 1919 年德文原文；文档里的中文是据此英译转写}
\source{16}{“权力趋于腐败，绝对权力绝对地腐败”}{Lord Acton 致 Mandell Creighton 主教信，1887 年 4 月 5 日；Hanover College 历史文本库节选页 \url{https://history.hanover.edu/courses/excerpts/165acton.html}，2026 年 9 月 18 日访问}{Power tends to corrupt and absolute power corrupts absolutely.}{2026-09-18}{已核实}{反方论点一}{该页注明文本转自 Online Library of Liberty，引句在第 7 段；紧接的下一句是 Great men are almost always bad men}
\source{17}{《民法典》第183、184条原文}{《中华人民共和国民法典》（2020 年 5 月 28 日第十三届全国人民代表大会第三次会议通过）第一百八十三条、第一百八十四条；全文见中华人民共和国国防部官网转载新华社通稿 \url{http://www.mod.gov.cn/gfbw/qwfb/yw_214049/4866121.html}}{因自愿实施紧急救助行为造成受助人损害的，救助人不承担民事责任}{2026-09-18}{已核实}{正方论点三}{第一百八十三条原文见同页同段：因保护他人民事权益使自己受到损害的，由侵权人承担民事责任，受益人可以给予适当补偿；全国人大国家法律法规数据库 flk.npc.gov.cn 本次未返回条文页}
\source{18}{院外心脏骤停存活出院率1.2\%、神经功能良好0.8\%、旁观者CPR 17.0\%、AED使用<0.1\%}{《中国心脏骤停与心肺复苏报告（2022年版）》，《中国循环杂志》2023年38卷10期1005–1017页}{}{}{有把握}{正方论点一}{本次只打开了健康界与澎湃新闻的转述，未打开期刊原文。赛前请队员在知网或万方检索该文核对四个数字}
\source{19}{义警行为的定义：法律程序之外对违法行为的预防、调查或惩罚}{Regina Bateson (2021), “The Politics of Vigilantism”, \emph{Comparative Political Studies} 54(6):923–955}{}{}{有把握}{双方定义调研}{SAGE 页面本次返回 403。赛前可在校内数据库或 ResearchGate 核对原文定义句}
\source{20}{安全困境：一方为自身安全采取的手段会降低他人安全}{Robert Jervis (1978), “Cooperation Under the Security Dilemma”, \emph{World Politics} 30(2):167–214}{}{}{有把握}{反方论点二}{在线 PDF 为扫描件无法提取文字。赛前用校内数据库核对第167–169页的定义段}
\source{21}{公共品供给不足与搭便车}{Mancur Olson (1965), \emph{The Logic of Collective Action}}{}{}{有把握}{正方论点一}{本次未查原文。核对第一章关于集团规模与公共品供给的论述}
\source{22}{个人化独裁比其他政体更容易发动战争}{Geddes, Wright \& Frantz (2018), \emph{How Dictatorships Work}, Cambridge University Press}{}{}{有把握}{反方论点三备用}{只见二手转述。赛前核对书中关于 personalist regimes 与对外冲突的章节}
\source{23}{中国每年院外心脏骤停约54.4万例}{该数字通行引用源自2005–2006年北京、广州、乌鲁木齐、阳泉四城市调查的全国外推}{}{}{待核实}{正方论点一备用}{查《中国心血管健康与疾病报告》（国家心血管病中心）历年版本，确认原始研究与外推方法。未核实前不要在场上说出这个数字}
\source{24}{Nelson \& Norton (2005) 启动效应是否有独立重复实验}{}{}{}{待核实}{双方论点三的攻防}{查 Open Science Collaboration (2015) 与 Many Labs 系列是否覆盖该范式。若对方质疑可重复性，需要这条来回应}
\source{25}{地震废墟埋压者存活率随时间衰减的具体曲线（“黄金72小时”）}{}{}{}{待核实}{正方论点一备用}{查 Macintyre, Barbera \& Smith (2006), \emph{Disaster Medicine and Public Health Preparedness}；或中国地震局历次救援统计。未核实前只讲机制不报数字}
\end{sourcelist}
\section{十二、备赛建议}
\begin{fields}
\begin{timeline}[时间表]
\when{距比赛 7 天}{全队背下第三节的中立定义与第四节的中立判准，逐字背，不要改词。四个人当场互相抽查。}
\when{距比赛 5 天}{正一、反一稿子成型并计时朗读；四辩开始背结辩；二辩、三辩把问题链背熟到不用看纸。}
\when{距比赛 3 天}{第一次模辩，重点检验下面三件事。}
\when{距比赛 1 天}{根据模辩复盘调整，把速查压到两页以内打印，第十一节里的三条【待核实】分工查完，查不到的当场删掉。}
\end{timeline}
\field{模辩重点检验的三件事}{
\begin{enumerate}[leftmargin=1.8em,itemsep=2pt,topsep=2pt]
\item 正方被问“谁来监督他”时，四个人的回答是不是同一套话（正确答案：承认不受问责，然后立刻回到判准问净变化，再反问“你说的灾难是和哪个世界比”）。
\item 反方被问“你凭什么假设他是坏人”时，会不会真的滑到“他迟早作恶”（正确答案：我方从不假设他是坏人，我方的问题是无法验证、也无法改回来）。
\item 双方有没有人在场上引用漫画情节当论据。有一个人说了，赛后全队重申一次证据纪律。
\end{enumerate}
}
\begin{checklist}[每位队员赛前要背下来的]
\item 中立定义三条（能力超常、主动介入、不受既有授权与问责）
\item 中立判准一句话（安全加自主，长期净变化）
\item 我方三个论点的标签
\item 我方主动坦白项（口径表最后一行），以及说出它的那句话
\item 对方最强的一个点，以及我方的一句话回应
\end{checklist}
\end{fields}
\end{document}
